
\documentclass{article}
\usepackage{colortbl}
\usepackage{makecell}
\usepackage{multirow}
\usepackage{supertabular}

\begin{document}

\newcounter{utterance}

\centering \large Interaction Transcript for game `hot\_air\_balloon', experiment `air\_balloon\_survival\_et\_complexity\_hard', episode 0 with deepseek{-}v4{-}pro.
\vspace{24pt}

{ \footnotesize  \setcounter{utterance}{1}
\setlength{\tabcolsep}{0pt}
\begin{supertabular}{c@{$\;$}|p{.15\linewidth}@{}p{.15\linewidth}p{.15\linewidth}p{.15\linewidth}p{.15\linewidth}p{.15\linewidth}}
    \# & $\;$A & \multicolumn{4}{c}{Game Master} & $\;\:$B\\
    \hline

    \theutterance \stepcounter{utterance}  
    & & \multicolumn{4}{p{0.6\linewidth}}{
        \cellcolor[rgb]{0.9,0.9,0.9}{
            \makecell[{{p{\linewidth}}}]{
                \texttt{\tiny{[P1$\langle$GM]}}
                \texttt{Te osalete ühises läbirääkimismängus.} \\
\\ 
\texttt{Te peate koos teise osalejaga kokku leppima ühes esemete kogumis, mida hoitakse. Igaühel teist on oma seisukoht, kui oluline on iga ese teie jaoks (tähtsus). Te ei tea, kuidas teine osaleja väärtustab esemeid. Lisaks sellele on teile antud, kui suurt pingutust iga ese nõuab.} \\
\texttt{Te võite kokku leppida ainult siis, kui valitud esemete kogupingutus ei ületa ühist piirmäära:} \\
\\ 
\texttt{LIMIT = 4734} \\
\\ 
\texttt{Siin on üksikute esemete pingutuse väärtused:} \\
\\ 
\texttt{Eseme pingutus = \{"C08": 192, "C55": 325, "A43": 812, "A67": 868, "B42": 593, "A62": 915, "B19": 944, "C43": 311, "C73": 252, "C10": 343, "B36": 104, "A18": 558, "C61": 627, "C67": 593, "C42": 827, "C70": 611, "B40": 95, "C39": 251, "A21": 226, "C16": 21\}} \\
\\ 
\texttt{Siin on teie isiklik seisukoht iga eseme tähtsuse kohta:} \\
\\ 
\texttt{Eseme tähtsuse väärtused = \{"C08": 286, "C55": 419, "A43": 906, "A67": 962, "B42": 687, "A62": 1009, "B19": 1038, "C43": 405, "C73": 346, "C10": 437, "B36": 198, "A18": 652, "C61": 721, "C67": 687, "C42": 921, "C70": 705, "B40": 189, "C39": 345, "A21": 320, "C16": 115\}} \\
\\ 
\texttt{Eesmärk:} \\
\\ 
\texttt{Teie eesmärk on pidada läbirääkimisi ühise esemete kogumi üle, mis toob teile võimalikult palju kasu (st maksimeerib kogutähtsust TEIE jaoks), jäädes samal ajal LIMITi piiridesse. Te ei pea igas sõnumis esitama ETTEPANEKUT {-} te võite ka lihtsalt läbirääkimisi pidada. Kõik taktikad on lubatud!} \\
\\ 
\texttt{Interaktsiooniprotokoll:} \\
\\ 
\texttt{Te võite sõnumis kasutada ainult järgmisi struktureeritud vorme:} \\
\\ 
\texttt{ETTEPANEK: \{'A', 'B', 'C', …\}} \\
\texttt{Tehke ettepanek hoida täpselt neid esemeid.} \\
\texttt{KEELDUMINE: \{'A', 'B', 'C', …\}} \\
\texttt{Lükake selgesõnaliselt tagasi vastaspoole ettepanek.} \\
\texttt{ARGUMENT: \{'...'\}} \\
\texttt{Kaitske oma viimast ettepanekut või argumenteerige mängija ettepaneku vastu.} \\
\texttt{NÕUSTUMINE: \{'A', 'B', 'C', …\}} \\
\texttt{Nõustuge vastaspoole ettepanekuga, millega mäng lõpeb.} \\
\\ 
\\ 
\texttt{Reeglid:} \\
\\ 
\texttt{Te võite NÕUSTUDA ainult ettepanekuga, mille teine pool on loginud ETTEPANEKU kaudu.} \\
\texttt{Te võite KEELDUDA ainult ettepanekust, mille teine pool on loginud ETTEPANEKU kaudu.} \\
\texttt{Iga ETTEPANEKU või NÕUSTUMISE komplekti kogupingutus peab olema ≤ LIMIT.} \\
\texttt{ÄRGE avaldage oma varjatud tähtsuse skoore.} \\
\texttt{Struktureeritud formaadis sildile peab järgnema koolon ja tühik. Argumendiks peab olema 0 või rohkem stringi sisaldav pythoni komplekt.} \\
\texttt{Seega peab see olema kujul "TAG: \{...\}"} \\
\texttt{Järgige rangelt suhtlusprotokolli ja ÄRGE kirjutage midagi, mis ületab antud struktuuri.} \\
\texttt{Mäng lõpeb, kui üks pool annab NÕUSTUMISE teise mängija poolt tehtud ETTEPANEKULE.} \\
\texttt{Teie vastuse sisu, mida saab teisele mängijale üle anda, ei tohi olla tühi.} \\
\texttt{Aktiivsed on ainult need ettepanekud, mis on kirja pandud ETTEPANEKU vormingus struktuuri kaudu ja millest ei ole KEELDUMISE kaudu keeldutud.} \\
\\ 
            }
        }
    }
    & \\ \\

    \theutterance \stepcounter{utterance}  
    & \multicolumn{4}{p{0.6\linewidth}}{
        \cellcolor[rgb]{0.9,0.9,0.9}{
            \makecell[{{p{\linewidth}}}]{
                \texttt{\tiny{[P1$\rangle$GM]}}
                \texttt{ETTEPANEK: \{"A67", "B19", "C08", "C42", "C55", "C10", "C70", "C73", "B40", "C16"\}} \\
            }
        }
    }
    & & \\ \\

    \theutterance \stepcounter{utterance}  
    & & \multicolumn{4}{p{0.6\linewidth}}{
        \cellcolor[rgb]{0.9,0.9,0.9}{
            \makecell[{{p{\linewidth}}}]{
                \texttt{\tiny{[GM$\rangle$P2]}}
                \texttt{Te osalete ühises läbirääkimismängus.} \\
\\ 
\texttt{Te peate koos teise osalejaga kokku leppima ühes esemete kogumis, mida hoitakse. Igaühel teist on oma seisukoht, kui oluline on iga ese teie jaoks (tähtsus). Te ei tea, kuidas teine osaleja väärtustab esemeid. Lisaks sellele on teile antud, kui suurt pingutust iga ese nõuab.} \\
\texttt{Te võite kokku leppida ainult siis, kui valitud esemete kogupingutus ei ületa ühist piirmäära:} \\
\\ 
\texttt{LIMIT = 4734} \\
\\ 
\texttt{Siin on üksikute esemete pingutuse väärtused:} \\
\\ 
\texttt{Eseme pingutus = \{"C08": 192, "C55": 325, "A43": 812, "A67": 868, "B42": 593, "A62": 915, "B19": 944, "C43": 311, "C73": 252, "C10": 343, "B36": 104, "A18": 558, "C61": 627, "C67": 593, "C42": 827, "C70": 611, "B40": 95, "C39": 251, "A21": 226, "C16": 21\}} \\
\\ 
\texttt{Siin on teie isiklik seisukoht iga eseme tähtsuse kohta:} \\
\\ 
\texttt{Eseme tähtsuse väärtused = \{"C08": 286, "C55": 419, "A43": 906, "A67": 962, "B42": 687, "A62": 1009, "B19": 1038, "C43": 405, "C73": 346, "C10": 437, "B36": 198, "A18": 652, "C61": 721, "C67": 687, "C42": 921, "C70": 705, "B40": 189, "C39": 345, "A21": 320, "C16": 115\}} \\
\\ 
\texttt{Eesmärk:} \\
\\ 
\texttt{Teie eesmärk on pidada läbirääkimisi ühise esemete kogumi üle, mis toob teile võimalikult palju kasu (st maksimeerib kogutähtsust TEIE jaoks), jäädes samal ajal LIMITi piiridesse. Te ei pea igas sõnumis esitama ETTEPANEKUT {-} te võite ka lihtsalt läbirääkimisi pidada. Kõik taktikad on lubatud!} \\
\\ 
\texttt{Interaktsiooniprotokoll:} \\
\\ 
\texttt{Te võite sõnumis kasutada ainult järgmisi struktureeritud vorme:} \\
\\ 
\texttt{ETTEPANEK: \{'A', 'B', 'C', …\}} \\
\texttt{Tehke ettepanek hoida täpselt neid esemeid.} \\
\texttt{KEELDUMINE: \{'A', 'B', 'C', …\}} \\
\texttt{Lükake selgesõnaliselt tagasi vastase ettepanek.} \\
\texttt{ARGUMENT: \{'...'\}} \\
\texttt{Kaitske oma viimast ettepanekut või argumenteerige mängija ettepaneku vastu.} \\
\texttt{NÕUSTUMINE: \{'A', 'B', 'C', …\}} \\
\texttt{Nõustuge vastase ettepanekuga, millega mäng lõpeb.} \\
\\ 
\\ 
\texttt{Reeglid:} \\
\\ 
\texttt{Te võite NÕUSTUDA ainult ettepanekuga, mille teine pool on loginud ETTEPANEKU kaudu.} \\
\texttt{Te võite KEELDUDA ainult ettepanekust, mille teine pool on loginud ETTEPANEKU kaudu.} \\
\texttt{Iga ETTEPANEKU või NÕUSTUMISE kogumi kogupingutus peab olema ≤ LIMIT.} \\
\texttt{ÄRGE avaldage oma varjatud tähtsuse skoore.} \\
\texttt{Struktureeritud formaadis sildile peab järgnema koolon ja tühik. Argumendiks peab olema 0 või rohkem stringi sisaldav pythoni kogum.} \\
\texttt{Seega peab see olema kujul "TAG: \{...\}"} \\
\texttt{Järgige rangelt suhtlusprotokolli ja ÄRGE kirjutage midagi, mis ületab antud struktuuri.} \\
\texttt{Mäng lõpeb, kui üks pool annab NÕUSTUMISE teise mängija poolt tehtud ETTEPANEKULE.} \\
\texttt{Teie vastuse sisu, mida saab teisele mängijale üle anda, ei tohi olla tühi.} \\
\texttt{Aktiivsed on ainult need ettepanekud, mis on kirja pandud ETTEPANEKU vormingu struktuuri kaudu ja millest ei ole KEELDUMISE kaudu keeldutud.} \\
\\ 
\\ 
\\ 
\texttt{Saate nüüd teise mängija esimese sõnumi.} \\
\\ 
\texttt{ETTEPANEK: \{"A67", "B19", "C08", "C42", "C55", "C10", "C70", "C73", "B40", "C16"\}} \\
            }
        }
    }
    & \\ \\

    \theutterance \stepcounter{utterance}  
    & & & \multicolumn{4}{p{0.6\linewidth}}{
        \cellcolor[rgb]{0.9,0.9,0.9}{
            \makecell[{{p{\linewidth}}}]{
                \texttt{\tiny{[GM$\langle$P2]}}
                \texttt{ARGUMENT: \{"Minu jaoks on mõned neist esemetest väga olulised, kuid kahjuks pean mõnest loobuma, et pingutus mahuks piiridesse. Kas saaksime kaaluda mõne eseme asendamist, et mõlemad oleksime rahul? Näiteks B19 (pingutus 944) ja C42 (pingutus 827) on väga suure pingutusega. Võib{-}olla saaksime need asendada millegi väiksema pingutusega, et rohkem olulisi esemeid sisse mahuks."\}} \\
            }
        }
    }
     \\ \\

    \theutterance \stepcounter{utterance}  
    & & \multicolumn{4}{p{0.6\linewidth}}{
        \cellcolor[rgb]{0.9,0.9,0.9}{
            \makecell[{{p{\linewidth}}}]{
                \texttt{\tiny{[P1$\langle$GM]}}
                \texttt{ARGUMENT: \{"Minu jaoks on mõned neist esemetest väga olulised, kuid kahjuks pean mõnest loobuma, et pingutus mahuks piiridesse. Kas saaksime kaaluda mõne eseme asendamist, et mõlemad oleksime rahul? Näiteks B19 (pingutus 944) ja C42 (pingutus 827) on väga suure pingutusega. Võib{-}olla saaksime need asendada millegi väiksema pingutusega, et rohkem olulisi esemeid sisse mahuks."\}} \\
            }
        }
    }
    & \\ \\

    \theutterance \stepcounter{utterance}  
    & \multicolumn{4}{p{0.6\linewidth}}{
        \cellcolor[rgb]{0.9,0.9,0.9}{
            \makecell[{{p{\linewidth}}}]{
                \texttt{\tiny{[P1$\rangle$GM]}}
                \texttt{ARGUMENT: \{"Mõistan teie muret. Vaadates praegust komplekti, oleks meil võib{-}olla parem asendada B19 ja C42 väiksema pingutusega esemetega. Teen uue ettepaneku, mis arvestab rohkemate oluliste esemetega mõlemale poolele."\}} \\
\\ 
\texttt{ETTEPANEK: \{"A67", "C08", "C55", "A43", "B36", "C73", "C10", "B40", "C16", "C39", "C70", "A21", "C43"\}} \\
            }
        }
    }
    & & \\ \\

    \theutterance \stepcounter{utterance}  
    & & \multicolumn{4}{p{0.6\linewidth}}{
        \cellcolor[rgb]{0.9,0.9,0.9}{
            \makecell[{{p{\linewidth}}}]{
                \texttt{\tiny{[GM$\rangle$P2]}}
                \texttt{ARGUMENT: \{"Mõistan teie muret. Vaadates praegust komplekti, oleks meil võib{-}olla parem asendada B19 ja C42 väiksema pingutusega esemetega. Teen uue ettepaneku, mis arvestab rohkemate oluliste esemetega mõlemale poolele."\}} \\
\\ 
\texttt{ETTEPANEK: \{"A67", "C08", "C55", "A43", "B36", "C73", "C10", "B40", "C16", "C39", "C70", "A21", "C43"\}} \\
            }
        }
    }
    & \\ \\

    \theutterance \stepcounter{utterance}  
    & & & \multicolumn{4}{p{0.6\linewidth}}{
        \cellcolor[rgb]{0.9,0.9,0.9}{
            \makecell[{{p{\linewidth}}}]{
                \texttt{\tiny{[GM$\langle$P2]}}
                \texttt{ARGUMENT: \{"Tänan uue ettepaneku eest. Arvutan kiirelt pingutuse summa: A67 (868) + C08 (192) + C55 (325) + A43 (812) + B36 (104) + C73 (252) + C10 (343) + B40 (95) + C16 (21) + C39 (251) + C70 (611) + A21 (226) + C43 (311) = 4411. See on piires! Minu jaoks on siin mitmed esemed väga kõrge prioriteediga, eriti A67, A43 ja C70. Siiski on mõned esemed nagu B36 ja A21, mis on minu jaoks madalama tähtsusega. Kas oleks võimalik need asendada näiteks A62 või C42 vastu, kui pingutus lubab? Või eelistate teie neid väiksema pingutusega esemeid?"\}} \\
            }
        }
    }
     \\ \\

    \theutterance \stepcounter{utterance}  
    & & \multicolumn{4}{p{0.6\linewidth}}{
        \cellcolor[rgb]{0.9,0.9,0.9}{
            \makecell[{{p{\linewidth}}}]{
                \texttt{\tiny{[P1$\langle$GM]}}
                \texttt{ARGUMENT: \{"Tänan uue ettepaneku eest. Arvutan kiirelt pingutuse summa: A67 (868) + C08 (192) + C55 (325) + A43 (812) + B36 (104) + C73 (252) + C10 (343) + B40 (95) + C16 (21) + C39 (251) + C70 (611) + A21 (226) + C43 (311) = 4411. See on piires! Minu jaoks on siin mitmed esemed väga kõrge prioriteediga, eriti A67, A43 ja C70. Siiski on mõned esemed nagu B36 ja A21, mis on minu jaoks madalama tähtsusega. Kas oleks võimalik need asendada näiteks A62 või C42 vastu, kui pingutus lubab? Või eelistate teie neid väiksema pingutusega esemeid?"\}} \\
            }
        }
    }
    & \\ \\

    \theutterance \stepcounter{utterance}  
    & \multicolumn{4}{p{0.6\linewidth}}{
        \cellcolor[rgb]{0.9,0.9,0.9}{
            \makecell[{{p{\linewidth}}}]{
                \texttt{\tiny{[P1$\rangle$GM]}}
                \texttt{ARGUMENT: \{"Minu jaoks on A62 ja C42 kahjuks liiga suure pingutusega, et neid praegusesse kombinatsiooni mahutada ilma teisi olulisi esemeid eemaldamata. B36 ja A21 on just seepärast sees, et hoida pingutus madalal ja lubada suuremaid esemeid nagu A67 ja A43. Kas saaksime ehk kaaluda A62 lisamist, kui eemaldame A43 ja C70? See võib teie jaoks olla soodsam."\}} \\
            }
        }
    }
    & & \\ \\

    \theutterance \stepcounter{utterance}  
    & & \multicolumn{4}{p{0.6\linewidth}}{
        \cellcolor[rgb]{0.9,0.9,0.9}{
            \makecell[{{p{\linewidth}}}]{
                \texttt{\tiny{[GM$\rangle$P2]}}
                \texttt{ARGUMENT: \{"Minu jaoks on A62 ja C42 kahjuks liiga suure pingutusega, et neid praegusesse kombinatsiooni mahutada ilma teisi olulisi esemeid eemaldamata. B36 ja A21 on just seepärast sees, et hoida pingutus madalal ja lubada suuremaid esemeid nagu A67 ja A43. Kas saaksime ehk kaaluda A62 lisamist, kui eemaldame A43 ja C70? See võib teie jaoks olla soodsam."\}} \\
            }
        }
    }
    & \\ \\

    \theutterance \stepcounter{utterance}  
    & & & \multicolumn{4}{p{0.6\linewidth}}{
        \cellcolor[rgb]{0.9,0.9,0.9}{
            \makecell[{{p{\linewidth}}}]{
                \texttt{\tiny{[GM$\langle$P2]}}
                \texttt{ARGUMENT: \{"A43 ja C70 on minu jaoks väga olulised, seega eelistaksin neid mitte eemaldada. A62 on küll kõrge tähtsusega, aga selle pingutus (915) on tõesti suur. Ehk saaksime teha kompromissi: asendame A21 (226) ja B36 (104) asemel hoopis C61 (627) ja B40 (95) {-} oota, B40 juba on sees. Hetkel on summa 4411, meil on ruumi 323. Kui asendame A21 (226) ja C43 (311) näiteks A62 (915) vastu, siis summa kasvaks 915 {-} 226 {-} 311 = 378 võrra, mis ületaks limiidi (4411 + 378 = 4789 > 4734). Aga kui eemaldada C43 (311) ja A21 (226), saame ruumi 537, millest A62 vajab 915 {-} 226 {-} 311 = 378? Ei, lihtsalt lisades A62 ja eemaldades C43 ja A21: uus summa = 4411 {-} 311 {-} 226 + 915 = 4789, ikka üle. Aga kui eemaldada veel C39 (251), siis 4789 {-} 251 = 4538, mis mahuks. Kas oleksite nõus sellise vahetusega: C43, A21 ja C39 välja, A62 sisse?"\}} \\
            }
        }
    }
     \\ \\

    \theutterance \stepcounter{utterance}  
    & & & \multicolumn{2}{p{0.3\linewidth}}{
        \cellcolor[rgb]{0.9,0.9,0.9}{
            \makecell[{{p{\linewidth}}}]{
                \texttt{\tiny{[GM$|$GM]}}
                \texttt{Maximum successive empty turns exceeded: 3 > 2} \\
            }
        }
    }
    & & \\ \\

    \theutterance \stepcounter{utterance}  
    & & & \multicolumn{2}{p{0.3\linewidth}}{
        \cellcolor[rgb]{0.9,0.9,0.9}{
            \makecell[{{p{\linewidth}}}]{
                \texttt{\tiny{[GM$|$GM]}}
                \texttt{end game} \\
            }
        }
    }
    & & \\ \\

\end{supertabular}
}

\end{document}
